\documentclass[11pt]{article}
\usepackage[margin=1in]{geometry}
\usepackage[T1]{fontenc}
\usepackage[utf8]{inputenc}
\usepackage{graphicx}
\usepackage{hyperref}
\usepackage{xurl}
\hypersetup{hidelinks,breaklinks=true}

\title{Adaptation Measures That Best Protect Biodiversity Hotspots and Nearby Communities Under Worsening Heat and Drought}
\author{Generated by Codex}
\date{April 22, 2026}

\begin{document}
\maketitle

\section{Question}
Which adaptation measures most protect biodiversity hotspots and nearby communities under worsening heat and drought?

\section{Overview}
The strongest cross-cutting evidence does not support a single stand-alone measure. It supports integrated packages that combine climate-smart area protection, landscape connectivity, protection of climate refugia, native ecosystem restoration and water-retention interventions, and governance that gives Indigenous Peoples and local communities real authority over land and resource decisions. Across reviews and syntheses, these measures are the most consistent at producing joint biodiversity and livelihood benefits under climate stress, whereas narrow or externally imposed measures more often underperform or create trade-offs.

\section{Main Findings}
First, protecting and connecting climate refugia appears to be one of the most robust hotspot-scale strategies. A review of recommendations for protected areas under climate change found five recurring priorities: ensure connectivity, protect climate refugia, favor a few large areas, protect places expected to become important in the future, and complement permanent reserves with temporary protection. The review argues that traditional conservation remains necessary but is insufficient without explicit climate adaptation design [1]. Country-level global analysis likewise found progress on refugia and abiotic diversity, but weaker progress on connectivity, indicating that reserve expansion often still misses the movement needs created by warming [2]. In the Mesoamerican hotspot, regional corridor analysis identified thousands of elevational climate-adaptation corridors, but many high-elevation refugia remained poorly protected, implying that hotspot protection needs to target both the corridor and the destination refuge [3].

Second, ecosystem-based adaptation and other nature-based measures have the clearest evidence for simultaneous biodiversity and community benefits when they rely on native ecosystems rather than simplified replacements. A systematic review of adaptation-focused nature-based solutions reported that 88\% of interventions with positive adaptation outcomes also improved ecosystem health, with a 67\% average increase in species richness, but noted trade-offs for some forest-management approaches and for novel or monoculture systems [4]. Reviews of drought-focused nature-based solutions point especially to wetland and floodplain restoration, rainwater harvesting, green infrastructure, and watershed-scale approaches as promising for buffering drought impacts, though large-scale quantitative evaluations remain limited [5]. Evidence from dryland watershed restoration in the Sky Islands hotspot suggests that rock detention structures adapted from Indigenous practice can increase infiltration, reduce peak runoff, retain sediment, improve vegetative cover, and support downstream livelihoods, making water-retention restoration particularly relevant where heat and drought reinforce erosion and flash flooding [6].

% Auto-generated by prepare_report_images.py
% Input this file where figures should appear in the review.

% Suggested section: Main Findings
\begin{figure}[htbp]
\centering
\includegraphics[width=0.92\linewidth]{literature_review_assets/images/01_central_america_corridor_map.png}
\caption{Regional study area showing lowland protected areas and high-elevation target areas used to map potential climate adaptation corridors across Central America. \emph{Source: McCullough et al. 2024, PLOS ONE}}
\label{fig:corridor-map}
\end{figure}

% Suggested section: Main Findings
\begin{figure}[htbp]
\centering
\includegraphics[width=0.92\linewidth]{literature_review_assets/images/02_central_america_corridor_characteristics.png}
\caption{Ecological characteristics and protection status of potential climate adaptation corridors by country, illustrating how corridor quality and protection vary across the hotspot region. \emph{Source: McCullough et al. 2024, PLOS ONE}}
\label{fig:corridor-characteristics}
\end{figure}


Third, restoration quality matters more than restoration counted as area alone. A global synthesis comparing restoration approaches found that native forests outperform compositionally simple plantations for biodiversity, water provisioning, and soil-erosion control, with younger plantations in drier regions performing especially poorly for non-wood outcomes [7]. A broader meta-analysis of conservation action also found large positive effects from habitat-loss reduction, restoration, invasive-species control, protected areas, and sustainable management, reinforcing that active conservation works when implemented at sufficient scale [8]. For heat- and drought-exposed hotspots, this implies that restoration should prioritize native composition, structural heterogeneity, and hydrological function rather than tree cover alone.

Fourth, community-led and Indigenous-led governance is among the strongest measures for protecting both biodiversity and nearby communities. A systematic review of 169 publications found that locally controlled conservation most often produced the strongest combined conservation and well-being outcomes, whereas externally controlled efforts that overrode customary institutions were frequently less effective ecologically and worse socially [9]. Related reviews of Indigenous and local knowledge show that adaptation portfolios commonly include land, soil, and water management, diversification, local governance, and traditional forecasting, and that these practices are most durable when rights, tenure, and institutions support them [10, 11, 12]. For hotspot regions where communities rely directly on forests, water, fisheries, or rain-fed agriculture, governance design is therefore not secondary to adaptation effectiveness; it is one of its main determinants.

Fifth, in fire-prone hotspots facing compound heat and drought, restoring resilient fire regimes is more protective than relying on suppression alone. Reviews for western North American forests support intentional management through prescribed burning, managed wildfire, and thinning paired with burning, especially where fire exclusion has increased vulnerability [13, 14]. Evidence from the Australian megafires shows biodiversity impacts were worse under extreme drought, outside protected areas, and where fire recurrence was high, supporting strategies that reduce repeat severe fire, defend wet ecosystems, and expand protection around vulnerable habitats [15]. This is a hotspot-specific measure rather than a universal one, but where heat and drought are amplifying fire risk, it belongs near the top of the adaptation list.

Sixth, measures that work best are usually bundled rather than isolated. Reviews of ecosystem-based and community-based adaptation repeatedly emphasize participatory governance, livelihood diversification, capacity building, multi-scalar planning, and long-term monitoring [16, 17, 18]. The practical implication is that hotspot adaptation is strongest when reserve design, restoration, water management, and local governance are implemented as a linked social-ecological program rather than as separate sector projects.

\section{Gaps or Limitations}
Most of the available evidence is indirect for the exact question of ``which measure most protects hotspots and nearby communities under worsening heat and drought.'' Much of the literature is based on reviews, conceptual syntheses, case studies, or observational evidence rather than head-to-head comparative trials. Evidence is also geographically uneven, with many syntheses leaning toward North America, Europe, forests, or coastal systems. Quantitative evidence on drought-specific nature-based adaptation remains relatively thin, and many studies assess biodiversity or social outcomes, but not both over long periods [4, 5, 19]. The literature also shows that some commonly promoted actions, especially monoculture plantations or top-down conservation, can look successful on paper while underdelivering on biodiversity or equity in practice [7, 9].

\section{Bottom Line}
The measures most consistently supported by the literature are: protect climate refugia and connect them across elevational or hydrological gradients; restore native ecosystems and watershed function rather than simple tree cover; give Indigenous Peoples and local communities real governance power, secure rights, and a central implementation role; and, in fire-prone systems, restore climate-resilient fire regimes instead of depending on suppression alone. If forced to prioritize, the best-supported package is climate-smart protected-area design plus native ecosystem and water-retention restoration, implemented through locally led governance. That combination appears to offer the strongest chance of protecting both biodiversity hotspots and the communities living alongside them as heat and drought intensify.

\section{References}
\begin{enumerate}
\item Ranius T, et al. Protected area designation and management in a world of climate change: A review of recommendations. Ambio. 2022. Available at: \url{https://consensus.app/papers/details/0b79634a19955c088ff3b086c59a1088/?utm_source=unknown}
\item Carrasco L, et al. Global progress in incorporating climate adaptation into land protection for biodiversity since Aichi targets. Global Change Biology. 2021. Available at: \url{https://consensus.app/papers/details/884997e7420e568e9968a56562b02255/?utm_source=unknown}
\item McCullough IM, et al. Mapping climate adaptation corridors for biodiversity---A regional-scale case study in Central America. PLOS ONE. 2024. Available at: \url{https://consensus.app/papers/details/9c6b0214be29562c8bc3851ce2544d86/?utm_source=unknown}
\item Key I, et al. Biodiversity outcomes of nature-based solutions for climate change adaptation: Characterising the evidence base. 2022. Available at: \url{https://consensus.app/papers/details/bff20cbfab8c558382b152aed889c197/?utm_source=unknown}
\item Yimer E A, et al. The underexposed nature-based solutions: A critical state-of-art review on drought mitigation. Journal of Environmental Management. 2024. Available at: \url{https://consensus.app/papers/details/fc07bef9e76557739d760983098ed690/?utm_source=unknown}
\item Gooden JL, et al. Dryland Watershed Restoration With Rock Detention Structures: A Nature-based Solution to Mitigate Drought, Erosion, Flooding, and Atmospheric Carbon. 2021. Available at: \url{https://consensus.app/papers/details/9d612f33b8615b3399b5fd9bfb53c699/?utm_source=unknown}
\item Hua F, et al. The biodiversity and ecosystem service contributions and trade-offs of forest restoration approaches. Science. 2022. Available at: \url{https://consensus.app/papers/details/f18afded537d574a90c279995eee7957/?utm_source=unknown}
\item Langhammer P, et al. The positive impact of conservation action. Science. 2024. Available at: \url{https://consensus.app/papers/details/6879a8a3b6e4514ab273e580624c9602/?utm_source=unknown}
\item Dawson N, et al. The role of Indigenous peoples and local communities in effective and equitable conservation. Ecology and Society. 2021. Available at: \url{https://consensus.app/papers/details/a2cb583e344e5bc98bc626c943ced53d/?utm_source=unknown}
\item Schlingmann A, et al. Global patterns of adaptation to climate change by Indigenous Peoples and local communities. A systematic review. Current Opinion in Environmental Sustainability. 2021. Available at: \url{https://consensus.app/papers/details/7178da6e27cf54fa8f8e0c54456f815e/?utm_source=unknown}
\item Dorji T, et al. Understanding How Indigenous Knowledge Contributes to Climate Change Adaptation and Resilience: A Systematic Literature Review. Environmental Management. 2024. Available at: \url{https://consensus.app/papers/details/cba760831d3353a7a3dae9011183fe4c/?utm_source=unknown}
\item Galappaththi E, et al. The sustainability assessment of Indigenous and local knowledge-based climate adaptation responses in agricultural and aquatic food systems. Current Opinion in Environmental Sustainability. 2023. Available at: \url{https://consensus.app/papers/details/e15ce8f8c9fb56229970ef617aaaab32/?utm_source=unknown}
\item Hessburg P, et al. Wildfire and climate change adaptation of western North American forests: a case for intentional management. Ecological Applications. 2021. Available at: \url{https://consensus.app/papers/details/a6bf2e62af5656e8b1550d5198699c04/?utm_source=unknown}
\item Prichard S, et al. Adapting western North American forests to climate change and wildfires: 10 common questions. Ecological Applications. 2021. Available at: \url{https://consensus.app/papers/details/2c152a5c52be5c0ba5a8432678174e4a/?utm_source=unknown}
\item Driscoll D, et al. Biodiversity impacts of the 2019--2020 Australian megafires. Nature. 2024. Available at: \url{https://consensus.app/papers/details/07f60966941b5f779e5c8da47e79924a/?utm_source=unknown}
\item Donatti C, et al. Indicators to measure the climate change adaptation outcomes of ecosystem-based adaptation. Climatic Change. 2019. Available at: \url{https://consensus.app/papers/details/e6e0b058cf49567eb665df513b31e2a4/?utm_source=unknown}
\item Selje T, et al. Community-Based Adaptation to Climate Change: Core Issues and Implications for Practical Implementations. Climate. 2024. Available at: \url{https://consensus.app/papers/details/418d9ff3b1d45400be8e28a5c02a7c1f/?utm_source=unknown}
\item McNamara K, et al. An assessment of community-based adaptation initiatives in the Pacific Islands. Nature Climate Change. 2020. Available at: \url{https://consensus.app/papers/details/1ced4e22c2d652c79fecabf1367c71da/?utm_source=unknown}
\item Morelli TL, et al. Conserving climate-change refugia: Insights from research and practice. Conservation Science and Practice. 2025. Available at: \url{https://consensus.app/papers/details/c80de269d9b35bc6ad59fd4c7ad14108/?utm_source=unknown}
\end{enumerate}

\end{document}
